%% Requires compilation with XeLaTeX or LuaLaTeX
\documentclass[10pt,xcolor={table,dvipsnames},t]{beamer}
\usetheme{UVic}

\usepackage{booktabs}
\usepackage{amsmath}
\usepackage{amssymb}

\title[PhD Proposal]{PhD Proposal}
\subtitle{Learning Automata Framework for Adaptive Kelp Forest Monitoring in British Columbia Coastal Waters}
\author{Ekaba Bisong}
\institute{Department of Computer Science\\University of Victoria\\[0.5em]Supervisor: Dr. Neil Ernst}
\date{January 2026}

\begin{document}

% Slide 1: Title
\begin{frame}
  \titlepage
\end{frame}

% Slide 2: Outline
\begin{frame}{Outline}
\tableofcontents
\end{frame}

% ============================================================================
\section{Introduction}
% ============================================================================

% Slide 3: Kelp Forest Importance
\begin{frame}{Kelp Forest Ecosystems}

\textbf{Ecological Significance:}
\begin{itemize}
  \item Among Earth's most productive ecosystems
  \item Support biodiversity comparable to tropical rainforests
  \item Provide critical habitat for marine species
  \item Carbon sequestration and coastal protection
\end{itemize}

\vspace{0.5em}
\textbf{Threats:}
\begin{itemize}
  \item Climate change (marine heatwaves)
  \item Urchin overgrazing
  \item Pollution and coastal development
  \item BC kelp forests declining at alarming rates
\end{itemize}

\end{frame}

% Slide 4: Monitoring Challenge
\begin{frame}{The Monitoring Challenge}

\textbf{Current Limitations:}
\begin{itemize}
  \item Traditional surveys: Labor-intensive, limited spatial coverage
  \item Satellite remote sensing: Scalable but faces challenges
  \begin{itemize}
    \item Variable environmental conditions (turbidity, SST)
    \item Seasonal and interannual changes
    \item Limited ground-truth labels
    \item Delayed validation data
  \end{itemize}
\end{itemize}

\vspace{0.5em}
\textbf{Key Insight:}
\begin{itemize}
  \item No single detection algorithm works best everywhere
  \item Need adaptive selection based on local conditions
  \item Must learn from sparse, delayed feedback
\end{itemize}

\end{frame}

% Slide 5: Problem Statement
\begin{frame}{Problem Statement}

\textbf{Core Challenge:}

How can we develop an automated kelp monitoring system that:

\begin{enumerate}
  \item Adapts to heterogeneous environmental conditions across BC's coastline?

  \item Learns effectively with limited labeled training data?

  \item Handles sparse and delayed ground-truth validation?

  \item Integrates Traditional Ecological Knowledge respectfully?
\end{enumerate}

\vspace{0.5em}
\textbf{Proposed Solution:} A Learning Automata framework for adaptive heuristic selection

\end{frame}

% ============================================================================
\section{Background}
% ============================================================================

% Slide 6: Learning Automata Overview
\begin{frame}{Learning Automata: Overview}

\textbf{Definition:} Adaptive decision-making agents that learn optimal actions through environment feedback

\vspace{0.5em}
\textbf{Key Components:}
\begin{itemize}
  \item \textbf{Actions} $\mathcal{A} = \{\alpha_1, \alpha_2, \ldots, \alpha_r\}$: Set of possible choices
  \item \textbf{Probability vector} $\mathbf{p}(n) = [p_1(n), \ldots, p_r(n)]$: Selection probabilities
  \item \textbf{Feedback} $\beta \in \{0, 1\}$: Reward (0) or penalty (1) from environment
  \item \textbf{Update rule}: Adjusts probabilities based on feedback
\end{itemize}

\vspace{0.5em}
\textbf{Goal:} Converge to optimal action that minimizes expected penalty

\end{frame}

% Slide 7: LA Update Schemes
\begin{frame}{Learning Automata: Update Schemes}

\textbf{Linear Reward-Inaction (L$_{\text{R-I}}$):}

On reward ($\beta = 0$) for action $\alpha_i$:
\begin{align}
p_i(n+1) &= p_i(n) + \lambda(1 - p_i(n)) \\
p_j(n+1) &= (1-\lambda)p_j(n), \quad j \neq i
\end{align}

On penalty ($\beta = 1$): No update (inaction)

\vspace{0.5em}
\textbf{Properties:}
\begin{itemize}
  \item $\lambda \in (0,1)$: Learning rate
  \item $\epsilon$-optimal in all stationary environments
  \item Robust to noisy feedback
  \item Well-suited for sparse reward scenarios
\end{itemize}

\end{frame}

% Slide 8: Non-Stationary Environments
\begin{frame}{Non-Stationary Environments}

\textbf{Markovian Switching Environment (MSE):}
\begin{itemize}
  \item Environment states form a Markov chain
  \item Unpredictable transitions between states
  \item Models interannual variability (climate anomalies)
\end{itemize}

\vspace{0.5em}
\textbf{Periodic Switching Environment (PSE):}
\begin{itemize}
  \item Cyclic state transitions (round-robin schedule)
  \item Predictable periodicity
  \item Models seasonal dynamics (kelp phenology)
\end{itemize}

\vspace{0.5em}
\textbf{Relevance:} Framework must handle both MSE-like (regime shifts) and PSE-like (seasonal) adaptation

\end{frame}

% ============================================================================
\section{Research Questions}
% ============================================================================

% Slide 9: Research Questions
\begin{frame}{Research Questions}

\begin{itemize}
  \item \textbf{RQ1 (Adaptation):} How effectively can an LA framework adapt to new geographic regions with limited labeled examples?

  \vspace{0.5em}
  \item \textbf{RQ2 (Heterogeneity):} Can context-aware heuristic selection improve detection accuracy across heterogeneous coastal environments?

  \vspace{0.5em}
  \item \textbf{RQ3 (Sparse Labels):} How can delayed and sparse ground-truth feedback be effectively incorporated into the learning process?

  \vspace{0.5em}
  \item \textbf{RQ4 (TEK Integration):} What protocols enable respectful and effective integration of Traditional Ecological Knowledge?
\end{itemize}

\end{frame}

% ============================================================================
\section{Proposed Methodology}
% ============================================================================

% Slide 10: Framework Overview
\begin{frame}{Framework Overview}

\textbf{Three Integrated Components:}

\begin{enumerate}
  \item \textbf{Component 1: Global Policy Learning}
  \begin{itemize}
    \item L$_{\text{R-I}}$ updates for heuristic selection
    \item Maintains global probability distribution over heuristics
  \end{itemize}

  \item \textbf{Component 2: Context-Aware Local Adaptation}
  \begin{itemize}
    \item Memory bank with context vectors
    \item k-NN retrieval for local probability estimation
    \item Global-local blending
  \end{itemize}

  \item \textbf{Component 3: Sparse Label Handling}
  \begin{itemize}
    \item Prediction buffer for asynchronous labels
    \item Temporal decay for stale feedback
    \item Delayed reward propagation
  \end{itemize}
\end{enumerate}

\end{frame}

% Slide 11: Component 1 - Global Policy
\begin{frame}{Component 1: Global Policy Learning}

\textbf{Heuristic Pool:} $\mathcal{H} = \{h_1, h_2, \ldots, h_r\}$

\textbf{Selection:} Sample heuristic $h_i$ with probability $p_i(n)$

\textbf{Reward Signal:}
\begin{equation}
\beta = \begin{cases}
0 & \text{if prediction matches ground truth} \\
1 & \text{otherwise}
\end{cases}
\end{equation}

\textbf{L$_{\text{R-I}}$ Update:} On reward, increase $p_i$; on penalty, no change

\end{frame}

% Slide 12: Component 1 - Heuristic Pool
\begin{frame}{Component 1: Heuristic Pool}

\textbf{Three-Tier Architecture:}

\begin{itemize}
  \item \textbf{Tier 1: Spectral Indices}
  \begin{itemize}
    \item NDVI, FAI, NDWI, Kelp Index
  \end{itemize}

  \item \textbf{Tier 2: Classical Machine Learning}
  \begin{itemize}
    \item Random Forest, SVM, XGBoost, Logistic Regression
  \end{itemize}

  \item \textbf{Tier 3: Deep Learning}
  \begin{itemize}
    \item U-Net, DeepLabV3+, ResNet, Vision Transformer
  \end{itemize}
\end{itemize}

\end{frame}

% Slide 12: Component 2 - Context Features
\begin{frame}{Component 2: Context Features}

\textbf{Context Vector Categories:}
\begin{itemize}
  \item \textbf{Spectral:} Mean, variance per band
  \item \textbf{Spatial:} Texture entropy, homogeneity
  \item \textbf{Temporal:} Day of year, time since observation
  \item \textbf{Environmental:} SST, turbidity, tide height
  \item \textbf{Geographic:} Lat/lon, distance to shore
\end{itemize}

\vspace{0.5em}
\textbf{Purpose:} Enable local adaptation based on observation conditions

\end{frame}

% Slide 13: Component 2 - Memory & Blending
\begin{frame}{Component 2: Memory \& Blending}

\textbf{Memory Bank:} $\mathcal{M} = \{(\mathbf{c}_j, \mathbf{p}_j, R_j)\}$

\textbf{k-NN Retrieval:} Find similar contexts using FAISS

\vspace{0.3em}
\textbf{Local Probability:}
\begin{equation}
\mathbf{p}_{\text{local}}(\mathbf{c}) = \frac{\sum_{j \in \mathcal{N}_k} w_j \cdot R_j \cdot \mathbf{p}_j}{\sum_{j \in \mathcal{N}_k} w_j \cdot R_j}
\end{equation}

\textbf{Global-Local Blending:}
\begin{equation}
\mathbf{p}_{\text{final}} = \alpha \cdot \mathbf{p}_{\text{global}} + (1-\alpha) \cdot \mathbf{p}_{\text{local}}
\end{equation}

\end{frame}

% Slide 13: Component 3 - Prediction Buffer
\begin{frame}{Component 3: Prediction Buffer}

\textbf{Challenge:} Ground-truth labels arrive asynchronously, often delayed

\vspace{0.5em}
\textbf{Buffer Structure:}
\begin{equation}
\mathcal{B} = \{(t_j, \mathbf{x}_j, \mathbf{c}_j, i_j, \hat{\mathbf{y}}_j)\}
\end{equation}

\vspace{0.3em}
\textbf{Components:}
\begin{itemize}
  \item $t_j$: Prediction timestamp
  \item $\mathbf{x}_j$: Input observation
  \item $\mathbf{c}_j$: Context vector
  \item $i_j$: Selected heuristic index
  \item $\hat{\mathbf{y}}_j$: Prediction output
\end{itemize}

\end{frame}

% Slide 14: Component 3 - Temporal Decay
\begin{frame}{Component 3: Temporal Decay}

\textbf{When Label Arrives:}
\begin{enumerate}
  \item Match label to buffered prediction
  \item Compute reward signal $\beta$
  \item Apply temporal decay: $\gamma^{(t_{\text{now}} - t_{\text{pred}})}$
  \item Update policy with weighted reward
\end{enumerate}

\vspace{0.5em}
\textbf{Decay Properties:}
\begin{itemize}
  \item Downweights stale feedback
  \item Reflects environmental changes since prediction
  \item $\gamma \in (0, 1)$ controls decay rate
\end{itemize}

\end{frame}

% Slide 14: RQ1 Conceptual Model
\begin{frame}{RQ1: Adaptation with Limited Labels}

\textbf{Hypothesis:} LA framework achieves comparable accuracy to supervised methods using 10--20\% of labeled data

\vspace{0.5em}
\textbf{Mechanism:}
\begin{itemize}
  \item Unsupervised heuristics provide initial predictions
  \item Sparse labels guide policy refinement
  \item L$_{\text{R-I}}$ concentrates probability on best-performing heuristics
\end{itemize}

\vspace{0.5em}
\textbf{Validation:}
\begin{itemize}
  \item Vary label availability: 5\%, 10\%, 20\%, 50\%, 100\%
  \item Compare: Fixed best heuristic, random selection, supervised CNN
  \item Measure: F1-score, precision, recall across label budgets
\end{itemize}

\end{frame}

% Slide 15: RQ2 Conceptual Model
\begin{frame}{RQ2: Heterogeneous Region Coverage}

\textbf{Hypothesis:} Context-aware selection outperforms global policy by adapting to local conditions

\vspace{0.5em}
\textbf{Mechanism:}
\begin{itemize}
  \item Different heuristics excel in different conditions
  \item Spectral indices: Good in clear water
  \item ML/DL models: Better with complex patterns
  \item Context features identify local conditions
\end{itemize}

\vspace{0.5em}
\textbf{Validation:}
\begin{itemize}
  \item Stratify by: Turbidity levels, SST ranges, geographic zones
  \item Compare: Global-only vs. context-aware selection
  \item Measure: Per-stratum accuracy, cross-region transfer
\end{itemize}

\end{frame}

% Slide 16: RQ3 Conceptual Model
\begin{frame}{RQ3: Sparse Label Learning}

\textbf{Hypothesis:} Buffering with temporal decay enables learning despite asynchronous feedback

\vspace{0.5em}
\textbf{Mechanism:}
\begin{itemize}
  \item Buffer stores predictions until labels arrive
  \item Temporal decay downweights stale feedback
  \item Policy updates reflect relevance of feedback
\end{itemize}

\vspace{0.5em}
\textbf{Validation:}
\begin{itemize}
  \item Simulate delay distributions: Days, weeks, months
  \item Compare: Immediate-only, delayed-as-immediate, proposed method
  \item Measure: Convergence rate, final accuracy, stability
\end{itemize}

\end{frame}

% Slide 17: RQ4 Conceptual Model
\begin{frame}{RQ4: TEK Integration}

\textbf{Hypothesis:} TEK enhances monitoring through historical baselines, site prioritization, and validation

\vspace{0.5em}
\textbf{Integration Points:}
\begin{itemize}
  \item Site importance weights $w_s$ for culturally significant areas
  \item Temporal priors from traditional seasonal calendars
  \item Validation labels from elder observations
\end{itemize}

\vspace{0.5em}
\textbf{Principles:}
\begin{itemize}
  \item Indigenous data sovereignty
  \item Community-defined protocols
  \item Benefit sharing (results back to community)
  \item Collaborative partnership, not extraction
\end{itemize}

\end{frame}

% ============================================================================
\section{Experimental Design}
% ============================================================================

% Slide 18: Dataset - Imagery
\begin{frame}{Dataset: Imagery \& Study Area}

\textbf{Study Area:} British Columbia coastline

\vspace{0.5em}
\textbf{Satellite Imagery:}
\begin{itemize}
  \item Sentinel-2 multispectral (10m resolution)
  \item Multiple years to capture temporal variation
  \item Cloud-free composites per season
  \item Bands: RGB, NIR, SWIR for kelp detection
\end{itemize}

\end{frame}

% Slide 19: Dataset - Ground Truth & Environmental
\begin{frame}{Dataset: Ground Truth \& Environmental}

\textbf{Ground Truth Sources:}
\begin{itemize}
  \item Existing kelp surveys and mapping products
  \item Drone validation where available
  \item TEK-derived observations (with permission)
\end{itemize}

\vspace{0.5em}
\textbf{Environmental Data:}
\begin{itemize}
  \item Sea surface temperature (SST)
  \item Turbidity estimates
  \item Tide and wave data
\end{itemize}

\end{frame}

% Slide 19: Evaluation Protocol
\begin{frame}{Evaluation Protocol}

\textbf{Protocol A: Temporal Adaptation}
\begin{itemize}
  \item Train on Year 1, test on Years 2--3
  \item Assess adaptation to temporal drift
\end{itemize}

\textbf{Protocol B: Spatial Generalization}
\begin{itemize}
  \item Train on Region A, test on Regions B, C
  \item Leave-one-region-out cross-validation
\end{itemize}

\textbf{Protocol C: Sparse Label Simulation}
\begin{itemize}
  \item Vary label arrival delays and densities
  \item Test robustness to feedback sparsity
\end{itemize}

\textbf{Metrics:} F1-score, Precision, Recall, IoU, Per-class accuracy

\end{frame}

% Slide 20: Ablation Studies - Components
\begin{frame}{Ablation Studies: Components}

\textbf{Component Ablations:}
\begin{itemize}
  \item Full framework vs. global-only (no context)
  \item Full framework vs. no temporal decay
  \item Full framework vs. immediate labels only
\end{itemize}

\vspace{0.5em}
\textbf{Heuristic Pool Analysis:}
\begin{itemize}
  \item Tier 1 only vs. Tiers 1+2 vs. Full pool
  \item Contribution of each heuristic type
\end{itemize}

\end{frame}

% Slide 21: Ablation Studies - Hyperparameters
\begin{frame}{Ablation Studies: Hyperparameters}

\textbf{Hyperparameter Sensitivity:}
\begin{itemize}
  \item Learning rate $\lambda \in \{0.01, 0.05, 0.1, 0.2\}$
  \item Number of neighbors $k \in \{3, 5, 10, 20\}$
  \item Decay factor $\gamma \in \{0.9, 0.95, 0.99\}$
  \item Blending coefficient $\alpha \in \{0.3, 0.5, 0.7\}$
\end{itemize}

\vspace{0.5em}
\textbf{Analysis:}
\begin{itemize}
  \item Grid search over parameter combinations
  \item Sensitivity curves for each parameter
  \item Interaction effects between parameters
\end{itemize}

\end{frame}

% ============================================================================
\section{Timeline and Contributions}
% ============================================================================

% Slide 21: Timeline 2026
\begin{frame}{Timeline: 2026}

\textbf{Research Foundation:}
\begin{itemize}
  \item January: Proposal defense
  \item May: M1 --- Dataset and baselines ready
  \item October: M2 --- Core framework functional
  \item November: Paper 1 (RQ1) submission
\end{itemize}

\end{frame}

% Slide 22: Timeline 2027
\begin{frame}{Timeline: 2027}

\textbf{Experiments \& Completion:}
\begin{itemize}
  \item January: M3 --- Temporal experiments complete
  \item February: Paper 2 (RQ2) submission
  \item March: M4, M5 --- Spatial and sparse label experiments
  \item April: Paper 3 (RQ3) submission
  \item May: M6 --- Extensions and TEK integration
  \item June: Paper 4 (RQ4) submission
  \item August: M7 --- Thesis draft complete
  \item October: M8 --- Thesis defense
\end{itemize}

\end{frame}

% Slide 23: Expected Contributions - Scientific & Methodological
\begin{frame}{Contributions: Scientific \& Methodological}

\textbf{Scientific Contributions:}
\begin{itemize}
  \item Novel LA framework for environmental monitoring
  \item Context-aware heuristic selection methodology
  \item Sparse label learning with temporal decay
\end{itemize}

\vspace{0.5em}
\textbf{Methodological Contributions:}
\begin{itemize}
  \item Integration of LA theory with remote sensing
  \item TEK integration protocols respecting data sovereignty
  \item Evaluation framework for adaptive monitoring systems
\end{itemize}

\end{frame}

% Slide 24: Expected Contributions - Practical Impact
\begin{frame}{Contributions: Practical Impact}

\textbf{Practical Applications:}
\begin{itemize}
  \item Scalable kelp monitoring for BC coastline
  \item Framework transferable to other ecosystems
  \item Decision support for conservation efforts
\end{itemize}

\vspace{0.5em}
\textbf{Target Users:}
\begin{itemize}
  \item Marine researchers and ecologists
  \item Conservation organizations
  \item Indigenous communities and resource managers
  \item Government environmental agencies
\end{itemize}

\end{frame}

% Slide 23: Summary
\begin{frame}{Summary}

\begin{itemize}
  \item \textbf{Problem:} Adaptive kelp monitoring with limited labels and heterogeneous conditions

  \vspace{0.3em}
  \item \textbf{Solution:} Learning Automata framework with three components
  \begin{itemize}
    \item Global policy learning
    \item Context-aware local adaptation
    \item Sparse label handling
  \end{itemize}

  \vspace{0.3em}
  \item \textbf{Validation:} Four research questions with rigorous experimental design

  \vspace{0.3em}
  \item \textbf{Timeline:} October 2027 completion with four publications

  \vspace{0.3em}
  \item \textbf{Impact:} Scalable monitoring supporting kelp conservation in BC
\end{itemize}

\end{frame}

% Slide 24: Questions
\begin{frame}{}

\vspace{2em}
\centering
{\Huge Thank You}

\vspace{1.5em}
{\Large Questions?}

\vspace{2em}
\normalsize
Ekaba Bisong\\
ebisong@uvic.ca

\end{frame}

\end{document}
